\documentclass[12pt,a4paper]{article}
\usepackage[T1]{fontenc}
\usepackage[utf8]{inputenc}
\usepackage{lmodern}
\usepackage{amsmath,amssymb,microtype}
\usepackage[margin=25mm,headheight=16pt]{geometry}
\usepackage{fancyhdr,enumitem}
\usepackage[hidelinks,pdfauthor={Marco Vinicio Gutiérrez Montenegro},pdftitle={Determinantes: actividad final de la Semana 12}]{hyperref}
\pagestyle{fancy}
\fancyhf{}
\fancyhead[L]{\small MA1103 · Cálculo y Álgebra Lineal}
\fancyhead[R]{\small Semana 12}
\fancyfoot[L]{\small Escuela de Matemática · TEC}
\fancyfoot[R]{\thepage}
\renewcommand{\headrulewidth}{0.4pt}
\setlength{\parindent}{0pt}
\setlength{\parskip}{7pt}
\setlist[enumerate]{leftmargin=*,itemsep=4pt,topsep=4pt}
\newcommand{\problema}[2]{\section*{Problema #1. #2}}
\begin{document}
\begin{center}
  {\small Instituto Tecnológico de Costa Rica\\Escuela de Matemática}\par
  \vspace{8pt}
  {\LARGE\bfseries Determinantes}\par
  {\large Actividad final · Estudio asincrónico de la Semana 12}\par
  {\normalsize Marco Vinicio Gutiérrez Montenegro}
\end{center}
\vspace{4pt}
\hrule
\vspace{7pt}
\textbf{Entrega:} jueves 22 de octubre de 2026, a más tardar a las
\textbf{5:00 p.\,m.}, en la plataforma \textbf{TEC Digital} (hora de Costa Rica).

\textbf{Nombre:} \rule{85mm}{0.4pt}\quad\textbf{Grupo:} \rule{20mm}{0.4pt}

\textbf{Indicaciones.} Presenta un desarrollo ordenado y legible. Justifica las propiedades y operaciones utilizadas. Los valores numéricos sin procedimiento no sustituyen la argumentación. Esta práctica no incluye soluciones. El registro del recorrido es un documento separado.

\problema{1}{Demostración por propiedades}
Sin desarrollar por cofactores, demuestra que, para cualesquiera valores de $a,b,c,d$,
\[
 \begin{vmatrix}
  1&a&b\\
  2&2a&2b\\
  3&c&d
 \end{vmatrix}=0.
\]
Extrae el factor de la fila correspondiente y escribe una igualdad de determinantes que justifique la conclusión. Indica la propiedad utilizada.

\problema{2}{Cálculo por dos caminos}
Calcula el siguiente determinante de orden tres:
\[
 D=\begin{vmatrix}
  1&2&1\\
  2&5&3\\
  1&3&4
 \end{vmatrix}.
\]
Realiza dos procedimientos: uno usando linealidad en una fila o columna, y otro mediante operaciones elementales. Muestra una sola operación en cada paso, conserva los factores necesarios y comprueba que ambos procedimientos coincidan.

\problema{3}{Regla de Cramer}
Resuelve directamente el sistema
\[
 \begin{cases}
  x+y+z=2,\\
  2x-y+z=2,\\
  x+2y-z=5.
 \end{cases}
\]
Construye la matriz de coeficientes, calcula $D,D_x,D_y,D_z$ y determina $x,y,z$ mediante los cocientes de la regla de Cramer. Comprueba tu solución en las tres ecuaciones.

\vfill
\hrule
\smallskip
{\small\textbf{Recordatorio de entrega:} jueves 22 de octubre de 2026, hasta las 5:00 p.\,m., en TEC Digital.}
\end{document}
